The three corrections remove false alarms. They do so for three different reasons, and the reasons say how far each result extends.

\subsection{Why the rules behave as they do}

The binomial coverage test is oversized because it treats the nominal coverage as the coverage of the set. The coverage of a conformal set is a random variable that has the Beta law of \cref{prop:law}, and a test window that is long against the calibration window sees the fluctuation of that variable many times over: the variance of the count is the binomial variance times $(m+n+1)/(n+2)$ (\cref{prop:variance}), which is \VarRatioC{} at equal windows of 1,000 rows and \VarRatioD{} at a test window four times the calibration window. The size of the binomial test is therefore a function of the ratio $m/n$ and not of the calibration size alone, as \cref{fig:coverage-size} shows, and the exact size converges to the limit of \cref{prop:binomial-size}. The false alarms of the binomial test are not noise around a small error: at a ratio of 4 the test rejects more than a third of the exchangeable windows, and a validation team reading it on out-of-time windows that are longer than the calibration sample would conclude that most stable models have drifted.

The fixed AUC-drop cut is oversized at small evaluation samples for a different reason. The drop on the worst subsample has a positive mean and a standard deviation that falls with the evaluation sample (\cref{conj:scale}), so a fixed cut has a size that depends on the sample size: \ShRateFixedFiveAZ{} at 125 rows and \ShRateFixedFiveDZ{} at 1,000 rows for the cut 0.05, and the cut 0.10 changes from over-sized to nearly calibrated to slightly conservative over the same range. The permutation reference removes the dependence because its null law is drawn at the sample size of the evaluation set, and its size stays between \ShRatePermAZ{} and \ShRatePermCZ{}. The fixed cuts have more power, and the reason is size: a cut calibrated to the size 0.05 has the power of the reference to within the error.

The amber cut of the situation light is the least transferable of the three. The share of significant complainants has a null law that depends on the number of applicants and on the structure of the data (\cref{tab:situation-designs}), so the false-alarm rate of a cut is a property of a cut, a size of portfolio and a design of the data, and the proposition that bounds it from the mean alone is attained by laws that the data do not realise (\cref{prop:share}). The cut 0.20 has a size below the level on both plants at 400 and 1,000 applicants. That is a consequence of a conservative choice at 400 applicants and not a property of the cut: at 1,000 applicants on the threshold-bias plant the size is \SiRateCutTwentyLargeZ{} and the power at the gap 0.2 is \SiRateCutTwentyLargeB{}, against \SiRateCutTenLargeB{} for the cut 0.10 at the size \SiRateCutTenLargeZ{}.

\subsection{What the product should do}

Four practices follow from the measurements, and the corrected tests of the engine already follow the first three. The calibration size must be recorded with a conformal set, because the law needs it and a wrong value shifts the size (\cref{tab:coverage-ablation}); the engine stores \texttt{n\_cal} beside the covered count. The Beta-Binomial p-value sets the light and the binomial p-value is kept beside it, so that a reader who has used the binomial rule can see where the two differ; on the committed windows they differ in the light of one of five. The permutation reference needs at least \RedMinDraws{} draws to be able to set a red light, and the default 199 allows a p-value of 0.005. The fourth practice is not yet followed: the amber cut of the situation light is set from a measurement at one number of applicants and one design, and it should be re-measured, by \cref{eq:cut-rule}, on a plant that has the size and the neighbour structure of each portfolio to which the light is applied. At the committed sizes of \CmFairQMin{} to \CmFairQMax{} complainants, the null share is narrower than at 400 applicants, so the cut 0.20 is likely conservative there and the amber light is less sensitive than it could be at the same size; the measurement of that size is the next step.

\subsection{Negative and partial results}

Three of the predictions are refuted. The cut 0.20 loses power against the cut 0.10 at small gaps (H3c); the size measured at 400 applicants does not carry to 1,000 (H3d) or to a second design (H3e). The permutation reference has less power than the fixed cut 0.05 in every cell, by at least \ShPermPowerGapMin{}, and it has less than the z rule at 500 and 1,000 evaluation rows; its advantage is in the size at small samples and in the absence of a calibration, and not in power. The measured size of the permutation reference lies above the exact size of its p-value grid by about 0.01 at 39, 99 and 199 draws (\ShPermExcessZMin{} to \ShPermExcessZMax{} standard errors), so the reference is not exactly valid in the engine's construction, even if its size stays within the bound of the predictions. The Beta-Binomial correction does not lose power in the sense of \cref{conj:power}, but it is not powerful: at a coverage lost of 0.01 with equal windows of 2,000 rows it rejects \CovPowBBBTen{} of the replications at $b=1.10$, so a validation team that needs to detect a loss of 0.01 in coverage reliably needs longer windows.

\subsection{Position against the published methods}

The Beta law of the conditional coverage and the Beta-Binomial law of the empirical coverage on a validation window are published \citep{vovk2012,angelopoulos2021}, and the variance factor of \cref{prop:variance} is the published variance of that law. What the measurement adds is that the test a validation team would write from the binomial law, which is the natural one, has the sizes of \cref{tab:exact-size}, and that its apparent advantage in power is its size. The permutation reference is a Monte Carlo test with valid p-values \citep{phipson2010,davison1997}; what is added is its use on the worst-subsample drop of a resilience test \citep{sudjianto2023} with development rows as the null, and the measurement of what it costs. The situation-testing light uses the per-complainant test of Alvarez and Ruggieri \citep{alvarez2025} and adds nothing to it; the measurement concerns the portfolio-level share and its cut point. The corrections are ordinary statistical practice applied to a validation battery; the contribution is the measurement and the decision rules that it supports.
